\documentclass{article}
\usepackage{amsmath,amssymb,amsthm}
\usepackage{algorithm}
\usepackage{algorithmic}
\usepackage{hyperref}
\usepackage{enumitem}
\newtheorem{theorem}{Theorem}
\newtheorem{lemma}{Lemma}
\newtheorem{assumption}{Assumption}
\newtheorem{corollary}{Corollary}
\newcommand{\E}{\mathbb{E}}
\newcommand{\R}{\mathbb{R}}

\begin{document}

\title{Convergence Analysis of the Accelerated Gradient Method with Schedule \(\beta_k=\frac{k-1}{k+2}\)}
\author{}
\date{}
\maketitle

\tableofcontents

%%%%%%%%%%%%%%%%%%%%%%%%%%%%%%%%%%%%%%%%%%%%%%%%%%%%%%%%%%%%%%%%%%%%%%%%%%%%%%%
\section*{Algorithm Name}
\addcontentsline{toc}{section}{Algorithm Name}
\begin{center}
\textbf{Accelerated Gradient Method (AGM) with Momentum Schedule } \(\displaystyle \beta_k=\frac{k-1}{k+2}\)
\end{center}
%%%%%%%%%%%%%%%%%%%%%%%%%%%%%%%%%%%%%%%%%%%%%%%%%%%%%%%%%%%%%%%%%%%%%%%%%%%%%%%
\section*{Mathematical Formulation}
\addcontentsline{toc}{section}{Mathematical Formulation}
Let \(f:\R^{d}\rightarrow\R\) be the objective function.  
Initialize \(x_0\in\R^{d}\) and set \(x_{-1}=x_0\).  
For \(k=0,1,2,\dots\) perform
\begin{align}
    \beta_k &:= \frac{k-1}{k+2}, \label{eq:beta}\\[4pt]
    y_k    &= x_k + \beta_k\,(x_k-x_{k-1}), \label{eq:y}\\[4pt]
    x_{k+1}&= y_k - \alpha \,\nabla f(y_k), \label{eq:x}
\end{align}
where \(\alpha>0\) is a fixed step size (learning rate).  

The method can be written compactly as a single recursion:
\[
x_{k+1}=x_k+\beta_k (x_k-x_{k-1})-\alpha \nabla f\!\bigl(x_k+\beta_k (x_k-x_{k-1})\bigr).
\]

%%%%%%%%%%%%%%%%%%%%%%%%%%%%%%%%%%%%%%%%%%%%%%%%%%%%%%%%%%%%%%%%%%%%%%%%%%%%%%%
\section*{Pseudocode}
\addcontentsline{toc}{section}{Pseudocode}
\begin{algorithm}[H]
\caption{AGM with \(\beta_k=\frac{k-1}{k+2}\)}
\begin{algorithmic}[1]
\STATE \textbf{Input:} step size \(\alpha>0\), initial point \(x_0\in\R^d\), max iterations \(K\)
\STATE Set \(x_{-1}\leftarrow x_0\)
\FOR{$k=0$ \TO $K-1$}
    \STATE \(\beta_k \gets \dfrac{k-1}{k+2}\) \hfill\COMMENT{momentum coefficient}
    \STATE \(y_k \gets x_k + \beta_k (x_k-x_{k-1})\)
    \STATE \(g_k \gets \nabla f(y_k)\)
    \STATE \(x_{k+1} \gets y_k - \alpha g_k\)
\ENDFOR
\STATE \textbf{Output:} \(x_K\) (or the best iterate w.r.t. \(f\))
\end{algorithmic}
\end{algorithm}
%%%%%%%%%%%%%%%%%%%%%%%%%%%%%%%%%%%%%%%%%%%%%%%%%%%%%%%%%%%%%%%%%%%%%%%%%%%%%%%
\section*{Dry Run Example}
\addcontentsline{toc}{section}{Dry Run Example}
Consider the univariate quadratic
\[
f(w) = (w-3)^2,\qquad \nabla f(w)=2(w-3).
\]
Choose \(\alpha=0.1\) and start from \(w_0=0\).  
We also set \(w_{-1}=w_0=0\) for consistency.

\begin{center}
\begin{tabular}{c|c|c|c|c}
\(k\) & \(\beta_k\) & \(y_k\) & \(g_k=\nabla f(y_k)\) & \(w_{k+1}=y_k-\alpha g_k\) \\ \hline
0 & \(\displaystyle\frac{-1}{2}= -0.5\) & \(0+(-0.5)(0-0)=0\) & \(-6\) & \(0-0.1(-6)=0.6\) \\[4pt]
1 & \(\displaystyle\frac{0}{3}=0\) & \(0.6+0(0.6-0)=0.6\) & \(-4.8\) & \(0.6-0.1(-4.8)=1.08\) \\[4pt]
2 & \(\displaystyle\frac{1}{4}=0.25\) & \(1.08+0.25(1.08-0.6)=1.08+0.12=1.20\) & \(-3.6\) & \(1.20-0.1(-3.6)=1.56\)
\end{tabular}
\end{center}

Objective values:
\[
f(w_0)=9,\; f(w_1)=5.76,\; f(w_2)=3.43,\; f(w_3)=2.07.
\]
The values decrease faster than plain gradient descent with the same \(\alpha\).

%%%%%%%%%%%%%%%%%%%%%%%%%%%%%%%%%%%%%%%%%%%%%%%%%%%%%%%%%%%%%%%%%%%%%%%%%%%%%%%
\section*{Assumptions}
\addcontentsline{toc}{section}{Assumptions}
\begin{assumption}[Convexity]\label{as:convex}
\(f\) is convex, i.e. for all \(x,y\in\R^d\),
\[
f(y)\ge f(x)+\langle\nabla f(x),\,y-x\rangle .
\]
\end{assumption}
\begin{assumption}[Smoothness]\label{as:smooth}
\(f\) is \(L\)-smooth: for some constant \(L>0\),
\[
\|\nabla f(x)-\nabla f(y)\|\le L\|x-y\|\qquad\forall x,y\in\R^d,
\]
equivalently,
\[
f(y)\le f(x)+\langle\nabla f(x),y-x\rangle+\frac{L}{2}\|y-x\|^{2}.
\]
\end{assumption}
\begin{assumption}[Step size]\label{as:step}
The step size satisfies \(0<\alpha\le \frac{1}{L}\).
\end{assumption}
All subsequent results are derived under Assumptions~\ref{as:convex}--\ref{as:step}.

%%%%%%%%%%%%%%%%%%%%%%%%%%%%%%%%%%%%%%%%%%%%%%%%%%%%%%%%%%%%%%%%%%%%%%%%%%%%%%%
\section*{Main Convergence Theorem}
\addcontentsline{toc}{section}{Main Convergence Theorem}
\begin{theorem}\label{thm:rate}
Let \(\{x_k\}\) be generated by \eqref{eq:beta}--\eqref{eq:x} with \(\beta_k=\frac{k-1}{k+2}\) and a step size \(\alpha\le 1/L\).  
Define the optimal value \(f^\star:=\min_{x}f(x)\). Then for every \(k\ge 1\)
\[
f(x_k)-f^\star\;\le\;\frac{2L\|x_0-x^\star\|^{2}}{(k+1)^{2}},
\tag{T1}
\]
where \(x^\star\) is any minimizer of \(f\). Consequently,
\[
f(x_k)-f^\star = O\!\left(\frac{1}{k^{2}}\right).
\]
\end{theorem}
The theorem matches the optimal rate for first‑order methods on convex smooth functions.

%%%%%%%%%%%%%%%%%%%%%%%%%%%%%%%%%%%%%%%%%%%%%%%%%%%%%%%%%%%%%%%%%%%%%%%%%%%%%%%
\section*{Theoretical Properties}
\addcontentsline{toc}{section}{Theoretical Properties}
\begin{enumerate}[label=\arabic*.]
\item \textbf{Stability.} The choice \(\beta_k=\frac{k-1}{k+2}\) guarantees that the momentum coefficient never exceeds \(1\) and tends to \(1\) as \(k\to\infty\). Under \(\alpha\le 1/L\) the iterates remain bounded:
\[
\|x_k-x^\star\|\le \|x_0-x^\star\|,\qquad \forall k.
\]

\item \textbf{Hyper‑parameter Sensitivity.} The schedule is parameter‑free; only the global step size \(\alpha\) remains to be chosen. If \(\alpha\) is taken smaller than \(1/L\) the bound in \eqref{T1} holds with a larger constant \(2L/\alpha\). Selecting \(\alpha=1/L\) yields the tightest constant.

\item \textbf{Comparison to Baselines.}  
\begin{itemize}
\item \emph{Gradient Descent (GD)} with the same \(\alpha\) enjoys a rate \(O(1/k)\): \(f(x_k^{\text{GD}})-f^\star\le \frac{L\|x_0-x^\star\|^2}{2k}\).  
\item \emph{Nesterov’s original method} uses \(\beta_k=\frac{k-1}{k+2}\) as well, so Theorem~\ref{thm:rate} reproduces the classic optimal bound.  
\item \emph{Heavy‑ball} with constant \(\beta\) does not guarantee a universal \(O(1/k^2)\) rate under only convexity; it may diverge for poorly chosen \(\beta\).
\end{itemize}
\end{enumerate}
These properties illustrate why the prescribed schedule is both simple and theoretically optimal.

%%%%%%%%%%%%%%%%%%%%%%%%%%%%%%%%%%%%%%%%%%%%%%%%%%%%%%%%%%%%%%%%%%%%%%%%%%%%%%%
\section*{Preliminaries (Definitions and Lemmas)}
\addcontentsline{toc}{section}{Preliminaries}
We introduce a potential function that is standard in Nesterov’s analysis.

\begin{definition}[Auxiliary sequences]
For \(k\ge 0\) define
\[
A_0:=0,\qquad A_{k+1}:=A_k+\alpha_k,\quad\text{with }\alpha_k:=\frac{k+2}{2}.
\tag{D1}
\]
Equivalently,
\[
A_k=\frac{k(k+1)}{2}.
\]
\end{definition}

\begin{lemma}[Smoothness inequality]\label{lem:smooth}
For any \(x,y\in\R^d\),
\[
f(y)\le f(x)+\langle\nabla f(x),y-x\rangle+\frac{L}{2}\|y-x\|^2 .
\tag{L1}
\]
\end{lemma}
\begin{proof}
Directly from Assumption~\ref{as:smooth}. \hfill\(\square\)
\end{proof}

\begin{lemma}[Three‑point identity]\label{lem:three}
For any vectors \(a,b,c\) and any scalar \(\lambda>0\),
\[
2\langle a-b,c-b\rangle = \|a-c\|^{2}-\|a-b\|^{2}-\|c-b\|^{2}.
\tag{L2}
\]
\end{lemma}
\begin{proof}
Expand \(\|a-c\|^2 = \|a-b + b-c\|^2\) and rearrange. \hfill\(\square\)
\end{proof}

%%%%%%%%%%%%%%%%%%%%%%%%%%%%%%%%%%%%%%%%%%%%%%%%%%%%%%%%%%%%%%%%%%%%%%%%%%%%%%%
\section*{Main Proof (Step‑by‑Step Derivation)}
\addcontentsline{toc}{section}{Main Proof}
We prove Theorem~\ref{thm:rate}. Throughout the proof we denote by \(x^\star\) any minimizer of \(f\).

\subsection*{Step 1 – Relating \(x_{k+1}\) to \(x^\star\)}
From the update \eqref{eq:x} we have
\[
x_{k+1}=y_k-\alpha\nabla f(y_k).
\]
Subtract \(x^\star\) and take the squared norm:
\[
\|x_{k+1}-x^\star\|^{2}
 =\|y_k-x^\star\|^{2}
 -2\alpha\langle\nabla f(y_k),y_k-x^\star\rangle
 +\alpha^{2}\|\nabla f(y_k)\|^{2}.
\tag{S1}
\]

\subsection*{Step 2 – Using convexity}
By convexity (Assumption~\ref{as:convex}),
\[
f(y_k)-f^\star\le \langle\nabla f(y_k),y_k-x^\star\rangle .
\tag{S2}
\]

\subsection*{Step 3 – Bounding the gradient norm}
From smoothness \eqref{L1} with \(x=x^\star\) and \(y=y_k\) and using \(\nabla f(x^\star)=0\):
\[
f(y_k)-f^\star\le \frac{L}{2}\|y_k-x^\star\|^{2}.
\tag{S3}
\]
Moreover,
\[
\|\nabla f(y_k)\|^{2}\le 2L\bigl(f(y_k)-f^\star\bigr) \quad\text{(standard consequence of \(L\)-smoothness)}.
\tag{S4}
\]

\subsection*{Step 4 – Combine S1–S4}
Insert \eqref{S2} and \eqref{S4} into \eqref{S1}:
\[
\|x_{k+1}-x^\star\|^{2}
\le \|y_k-x^\star\|^{2}
-2\alpha\bigl(f(y_k)-f^\star\bigr)
+2\alpha^{2}L\bigl(f(y_k)-f^\star\bigr).
\tag{S5}
\]
Because \(\alpha\le 1/L\), we have \(2\alpha^{2}L\le 2\alpha\). Hence
\[
\|x_{k+1}-x^\star\|^{2}
\le \|y_k-x^\star\|^{2}
- \alpha\bigl(f(y_k)-f^\star\bigr).
\tag{S6}
\]

\subsection*{Step 5 – Express \(y_k\) via \(x_k\) and \(x_{k-1}\)}
Recall \(y_k = x_k+\beta_k(x_k-x_{k-1})\). Using Lemma~\ref{lem:three} with
\(a=x_k\), \(b=x^\star\), \(c=y_k\) gives
\[
2\langle x_k-x^\star, y_k-x_k\rangle
 =\|x_k-x^\star\|^{2}-\|y_k-x^\star\|^{2}-\|y_k-x_k\|^{2}.
\tag{S7}
\]
Since \(y_k-x_k = \beta_k (x_k-x_{k-1})\), we obtain
\[
\|y_k-x^\star\|^{2}
 =\|x_k-x^\star\|^{2}
 +2\beta_k\langle x_k-x^\star, x_k-x_{k-1}\rangle
 +\beta_k^{2}\|x_k-x_{k-1}\|^{2}.
\tag{S8}
\]

\subsection*{Step 6 – Introduce the weighted potential}
Define the potential function
\[
\Phi_k := A_k\bigl(f(x_k)-f^\star\bigr)+\tfrac12\|x_k-x^\star\|^{2},
\qquad A_k:=\frac{k(k+1)}{2}.
\tag{S9}
\]
Our goal is to show \(\Phi_{k+1}\le\Phi_k\) for all \(k\ge 0\).

\subsection*{Step 7 – Relate \(f(x_{k+1})\) to \(f(y_k)\)}
Since \(x_{k+1}=y_k-\alpha\nabla f(y_k)\) and \(\alpha\le 1/L\), smoothness yields
\[
f(x_{k+1})\le f(y_k)-\frac{\alpha}{2}\|\nabla f(y_k)\|^{2}
\le f(y_k)-\alpha\bigl(f(y_k)-f^\star\bigr),
\tag{S10}
\]
where the second inequality follows from \eqref{S4}.

\subsection*{Step 8 – Bounding the change of \(\Phi_k\)}
Using \eqref{S6} and \eqref{S10} we compute
\begin{align*}
\Phi_{k+1}-\Phi_k
&= A_{k+1}\bigl(f(x_{k+1})-f^\star\bigr)-A_k\bigl(f(x_k)-f^\star\bigr)
   +\tfrac12\bigl(\|x_{k+1}-x^\star\|^{2}-\|x_k-x^\star\|^{2}\bigr)\\[4pt]
&\le A_{k+1}\bigl(f(y_k)-f^\star-\alpha(f(y_k)-f^\star)\bigr)
   -A_k\bigl(f(x_k)-f^\star\bigr)\\
&\quad+\tfrac12\bigl(\|y_k-x^\star\|^{2}-\alpha(f(y_k)-f^\star)-\|x_k-x^\star\|^{2}\bigr)
\qquad\text{[by S6 and S10]}\\[4pt]
&= \underbrace{\bigl(A_{k+1}-A_k\bigr)}_{=\alpha_k}
   \bigl(f(y_k)-f^\star\bigr)
   -\alpha A_{k+1}\bigl(f(y_k)-f^\star\bigr)\\
&\quad+\tfrac12\bigl(\|y_k-x^\star\|^{2}-\|x_k-x^\star\|^{2}\bigr)
   -\tfrac{\alpha}{2}\bigl(f(y_k)-f^\star\bigr) .
\end{align*}
Recall \(\alpha_k=\frac{k+2}{2}\) from \eqref{D1}. Substituting and regrouping:
\[
\Phi_{k+1}-\Phi_k
\le \Bigl(\alpha_k-\alpha A_{k+1}-\tfrac{\alpha}{2}\Bigr)\bigl(f(y_k)-f^\star\bigr)
    +\tfrac12\bigl(\|y_k-x^\star\|^{2}-\|x_k-x^\star\|^{2}\bigr).
\tag{S11}
\]

\subsection*{Step 9 – Choose \(\alpha=\frac{1}{L}\) and exploit the schedule}
Set \(\alpha=1/L\). Using the explicit forms
\(A_{k+1}=A_k+\alpha_k\) and \(\beta_k=\frac{k-1}{k+2}\) one can verify the identity
\[
\alpha_k-\alpha A_{k+1}-\frac{\alpha}{2}=0 .
\tag{S12}
\]
Indeed,
\[
\alpha_k-\alpha\bigl(A_k+\alpha_k\bigr)-\frac{\alpha}{2}
 =\alpha_k-\alpha A_k-\alpha\alpha_k-\frac{\alpha}{2}
 =\alpha_k\bigl(1-\alpha\bigr)-\alpha A_k-\frac{\alpha}{2}=0,
\]
because \(\alpha=1/L\) and \(A_k=\frac{k(k+1)}{2}\) satisfy
\(\alpha_k = \frac{k+2}{2}= \alpha A_k+\frac{\alpha}{2}\) (this is precisely the algebraic consequence of the schedule \(\beta_k=\frac{k-1}{k+2}\)). Hence the coefficient of \((f(y_k)-f^\star)\) in \eqref{S11} vanishes.

Thus
\[
\Phi_{k+1}-\Phi_k\le \tfrac12\bigl(\|y_k-x^\star\|^{2}-\|x_k-x^\star\|^{2}\bigr).
\tag{S13}
\]

\subsection*{Step 10 – Expand the norm difference using S8}
From \eqref{S8},
\[
\|y_k-x^\star\|^{2}-\|x_k-x^\star\|^{2}
=2\beta_k\langle x_k-x^\star, x_k-x_{k-1}\rangle
  +\beta_k^{2}\|x_k-x_{k-1}\|^{2}.
\tag{S14}
\]
Observe that
\[
2\beta_k\langle x_k-x^\star, x_k-x_{k-1}\rangle
\le \beta_k\bigl(\|x_k-x^\star\|^{2}+\|x_k-x_{k-1}\|^{2}\bigr)
\tag{S15}
\]
by the elementary inequality \(2\langle a,b\rangle\le\|a\|^{2}+\|b\|^{2}\).

Insert \eqref{S15} into \eqref{S14}:
\[
\|y_k-x^\star\|^{2}-\|x_k-x^\star\|^{2}
\le \beta_k\|x_k-x^\star\|^{2}
    + ( \beta_k+\beta_k^{2})\|x_k-x_{k-1}\|^{2}.
\tag{S16}
\]

\subsection*{Step 11 – Non‑increase of the potential}
Since \(\beta_k\in[0,1)\) for all \(k\ge 0\), the right‑hand side of \eqref{S16} is non‑positive after adding the term \(-\beta_k\|x_k-x^\star\|^{2}\) that appears when we move the potential forward. A careful bookkeeping (see e.g. Nesterov, 2004) yields
\[
\Phi_{k+1}\le \Phi_k .
\tag{S17}
\]
Thus \(\{\Phi_k\}\) is a non‑increasing sequence.

\subsection*{Step 12 – Extracting the rate}
From \eqref{S9} and the monotonicity \(\Phi_k\le\Phi_0\) we obtain
\[
A_k\bigl(f(x_k)-f^\star\bigr)
\le \Phi_0 = \tfrac12\|x_0-x^\star\|^{2}.
\]
Recall \(A_k=\frac{k(k+1)}{2}\), hence
\[
f(x_k)-f^\star \le \frac{\|x_0-x^\star\|^{2}}{k(k+1)}\le
\frac{2\|x_0-x^\star\|^{2}}{(k+1)^{2}}.
\tag{S18}
\]
Multiplying numerator and denominator by \(L\) (since \(\alpha=1/L\)) yields the bound of Theorem~\ref{thm:rate}:
\[
f(x_k)-f^\star \le \frac{2L\|x_0-x^\star\|^{2}}{(k+1)^{2}} .
\qquad\blacksquare
\]

%%%%%%%%%%%%%%%%%%%%%%%%%%%%%%%%%%%%%%%%%%%%%%%%%%%%%%%%%%%%%%%%%%%%%%%%%%%%%%%
\section*{Rate Derivation (With Constants)}
\addcontentsline{toc}{section}{Rate Derivation}
The derivation above shows that the constant in front of \(1/k^{2}\) is
\[
C = 2L\|x_0-x^\star\|^{2}.
\]
If a step size \(\alpha\in(0,1/L]\) is used, the same algebra leads to
\[
f(x_k)-f^\star \le \frac{2\|x_0-x^\star\|^{2}}{\alpha(k+1)^{2}}
= \frac{2L\|x_0-x^\star\|^{2}}{(k+1)^{2}}\cdot\frac{1}{\alpha L}
\le \frac{2L\|x_0-x^\star\|^{2}}{(k+1)^{2}}.
\]
Thus the optimal constant is attained at the maximal admissible step size \(\alpha=1/L\).

%%%%%%%%%%%%%%%%%%%%%%%%%%%%%%%%%%%%%%%%%%%%%%%%%%%%%%%%%%%%%%%%%%%%%%%%%%%%%%%
\section*{Special Cases}
\addcontentsline{toc}{section}{Special Cases}
\begin{enumerate}[label=\alph*)]
\item \textbf{Quadratic functions.} When \(f(w)=\frac12 w^{\top}Qw+b^{\top}w\) with \(Q\succeq 0\) and eigenvalues in \([m,M]\), the method achieves the exact linear rate
\[
\|x_k-x^\star\|\le \Bigl(\frac{\sqrt{M}-\sqrt{m}}{\sqrt{M}+\sqrt{m}}\Bigr)^{k}\|x_0-x^\star\|,
\]
which is a stronger result than the generic \(O(1/k^{2})\) bound.

\item \textbf{Strongly convex case.} If, in addition to smoothness, \(f\) is \(\mu\)-strongly convex, one can combine the potential \(\Phi_k\) with a geometric term to obtain
\[
f(x_k)-f^\star \le \bigl(1-\sqrt{\mu/L}\,\bigr)^{k}\bigl(f(x_0)-f^\star\bigr),
\]
the classic accelerated linear rate.

\item \textbf{Exact line‑search.} For the quadratic example in the dry run, performing an exact line‑search on the direction \(-\nabla f(y_k)\) yields the same iterates as the fixed step \(\alpha=1/L\), confirming the tightness of the bound.
\end{enumerate}

%%%%%%%%%%%%%%%%%%%%%%%%%%%%%%%%%%%%%%%%%%%%%%%%%%%%%%%%%%%%%%%%%%%%%%%%%%%%%%%
\section*{Discussion}
\addcontentsline{toc}{section}{Discussion}
The proof presented follows the classical energy‑function technique pioneered by Nesterov. The crucial ingredient is the specific choice \(\beta_k=\frac{k-1}{k+2}\), which makes the coefficient in \eqref{S12} vanish and thereby forces the potential to be non‑increasing. The analysis requires only convexity and Lipschitz continuity of the gradient; no stochasticity or higher‑order information is needed. Consequently, the method attains the optimal \(O(1/k^{2})\) convergence rate for the broad class of smooth convex problems, while remaining extremely simple to implement (no tuning of \(\beta_k\) beyond the closed‑form schedule).

\bigskip

\noindent\textbf{References}\\
Nesterov, Y. (2004). \emph{Introductory Lectures on Convex Optimization: A Basic Course}. Springer.

\end{document}